\section{Group or place? Solo flights where groups fly}\label{sec:zones}

Do groups fly where, and when, any pilot would progress straighter? If so, solo flights compared on
the same ground and at the same hours would match group time.

\paragraph{Reweighting.} Each piece gets a stratum from its middle, in four stratifications from coarse
to fine: zone, zone$\times$month, zone$\times$hour, site-day$\times$hour. The solo flights are then given
the mix of the group over the strata (direct standardisation):
\begin{equation}\label{eq:reweight}
  \MSD_{\Sl\mid\G}(\tau)=\frac{\sum_s n_{\G,s}\,\MSD_{\Sl,s}(\tau)}{\sum_s n_{\G,s}},
\end{equation}
with $n_{\G,s}$ the units of $\G$ in stratum $s$ (a flight counts once per stratum it crosses) and
$\MSD_{\Sl,s}$ the ensemble MSD of the solo units of $s$. In the statistic, $\Sl$ then flies where
and when $\G$ flies. $\G$ is taken on the strata it shares with $\Sl$, so the two curves have the same
support; the same is done for $\A$. ``As is'' is the comparison on the same strata without the
reweighting.

\begin{figure}
\centering
\includegraphics[width=\textwidth]{report_zones_lags.pdf}
\caption{$\Delta\Hloc(\tau)$ over every lag: $\G$ minus $\Sl$ (left) and $\G$ minus $\A$ (right), the
other population reweighted to the strata of $\G$, with 5--95\% bands; one line per stratification,
finer strata darker, grey without strata. The curves themselves are in Figure~\ref{fig:zones}.}
\label{fig:zones-diff}
\end{figure}

\begin{figure}
\centering
\includegraphics[width=\textwidth]{report_zones.pdf}
\caption{$\Delta\Hloc(1000~\mathrm{s})$ by stratification, with the share of $\G$ kept: $\G-\Sl$ and
$\G-\A$, the other population reweighted to the strata of $\G$ (filled) and as is on the same strata
(open). Values with the fitted $\Delta H$ in Table~\ref{tab:zones}.}
\label{fig:zones-forest}
\end{figure}

\paragraph{Results.} The gap opens above $\tau\approx150$~s and grows up to the longest lag
(Figure~\ref{fig:zones-diff}). Reweighting the solo flights to the strata of the groups lowers the
long-lag gap of the same strata little (Figure~\ref{fig:zones-forest}, filled against open points):
by \ZoneRedLocZone\% with the zones alone and by at most a fifth with the months or the hours,
although solo pilots do fly where groups fly (\ZoneKeptZone\% of $\G$ has a solo counterpart in its
cell). Against the same flights alone the gap is smaller, and the reweighting lowers it by a similar
share. At the same site-day and hour the reweighted $\G-\A$ (\ZoneDHlocASdh) differs from the
contrast of Section~\ref{sec:other} on the same strata (\DHlocThousandGroupSdh) only through the
weights, $n_{\G,s}$ here and the half harmonic mean there: the effect varies between strata, so its
average depends on how they are weighted.

\paragraph{Limits.} Cells are 25~km, placed by the middle of a piece; the terrain inside a cell is
not controlled, and a group may choose better air within it, which is then part of what we call the
group effect. Paragliders only.

\paragraph{What we keep.} Reweighting the solo flights to the zones, months and hours where groups fly
lowers the long-lag gap by at most a fifth, much less with the zone alone; the rest stays at the same
site, day and hour. Inside a zone, the route itself remains to be tested.
